% Includable mathematical note. Uses only standard amsmath/amssymb commands.
% Written for the pinned ProveIt baseline 6ca328039cafa1537d81ee6b67ee2dc54ffa5a4e.

\section{Uniform successor evaluation on zero-normal terms}
\label{sec:uniform-step-theory}

\subsection{Canonical syntax and what it costs}

Start with raw scalar terms
\[
 t ::= x\mid 0\mid S(t)\mid I(n,a;u.h),
\]
where $u$ binds only its occurrences in $h$. Define $N$ by structural
recursion, normalizing all immediate subterms and replacing
$I(0,a;u.h)$ by $a$ whenever the normalized counter is literally zero.
Only this zero rule is contracted. In particular,
$I(S0,a;u.h)$ is retained and is not definitionally $h[a/u]$.

The numerical terms of the new calculus are the fixed points of $N$.
Every displayed substitution means ordinary capture-free substitution
followed by $N$; write $h\langle t\rangle=N(h[t/u])$. Normalization applies
also under all term and formula binders. On formulas it changes numerical
terms only. Freshness is computed in the usual free-variable sense.

The elementary coherence equations are
\begin{align*}
 N(N(t))&=N(t),\\
 N(t[\sigma])&=N(N(t)[N\circ\sigma]),\\
 (t\langle\sigma\rangle)\langle\tau\rangle
 &=t\langle x\mapsto\sigma(x)\langle\tau\rangle\rangle.
\end{align*}
They follow by structural induction, with the binder cases using lifted
substitutions. Normalization commutes with renaming and introduces no
free variables. The contraction procedure and normalized substitution
belong to the trusted specification; deleting the zero computation from
the inference rules does not delete this implementation cost.

\subsection{One evaluation pattern suffices for equality and computation}

Use reflexivity and just the following evaluation inference:
\[
 \frac{\Gamma\vdash I(c,a;u.h)=t
       \qquad
       \Gamma\vdash I(Sc,a;u.h)=r}
      {\Gamma\vdash r=h\langle t\rangle}.
 \tag{Step}\label{rule:uniform-step}
\]
All terms in the display are interpreted by the canonical convention.
In particular, at $c=0$ the first source is simply $a$. The same
$c,a,h$ must occur in both sources. There is one case of this rule, no
comparison of a term with an arbitrary context and no search for an
occurrence to replace.

\paragraph{Successor computation.}
Take reflexivity as both premises of Step. Its conclusion is precisely
\[
 I(Sc,a;u.h)=h\langle I(c,a;u.h)\rangle.
 \tag{Succ}
\]
Zero computation is already the definition of the canonical constructor.

\paragraph{Equality pivot.}
For an arbitrary canonical term $s$, abbreviate
\[
 J_s=I(S0,s;v.v).
\]
Given $p:\Gamma\vdash s=t$ and $q:\Gamma\vdash s=r$, Step at
$c=0,a=s,h(v)=v$, with major premise $q$ and minor premise the
reflexivity of $J_s$, proves $J_s=r$. A second instance of the same
Step, now with premises $p$ and $J_s=r$, proves $r=t$. Thus
\[
 \frac{\Gamma\vdash s=t\qquad\Gamma\vdash s=r}
      {\Gamma\vdash r=t}
 \tag{Pivot}\label{rule:pivot}
\]
is a derived rule using two Steps and one auxiliary reflexivity.
This derivation is valid at every $s$, including literal zero.

\paragraph{Symmetry and transitivity.}
Pivot applied to $s=s$ and $s=t$ gives $t=s$. Given $s=t$ and $t=r$,
first obtain $t=s$ by symmetry and then apply Pivot with common source
$t$ to obtain $s=r$. Hence all equivalence-relation laws of equality
are available.

\paragraph{Congruence in every fixed term template.}
Fix a canonical body $h(v)$, allowing passive parameters and nested
iterators, and put $K=I(S0,s;v.h(v))$. Step with the reflexivities of
$s$ and $K$ proves $K=h\langle s\rangle$. Step with $s=t$ and the
reflexivity of $K$ proves $K=h\langle t\rangle$. Pivot gives
\[
 \Gamma\vdash h\langle s\rangle=h\langle t\rangle.
\]
Combining this with $h\langle s\rangle=r$ by Pivot gives the
one-template crossed transport rule
\[
 \frac{\Gamma\vdash s=t\qquad
       \Gamma\vdash h\langle s\rangle=r}
      {\Gamma\vdash r=h\langle t\rangle}.
 \tag{Tr}\label{rule:recovered-tr}
\]
Two instances give transport in an equation template. None of these
arguments uses equation induction, numerical case distinction,
successor injectivity, a logical connective, or a structural term
substitution theorem for whole derivations.

\paragraph{Reverse derivation.}
In the zero-normal native kernel, successor computation gives
$I(Sc,a;h)=h\langle I(c,a;h)\rangle$. The two premises of Step,
symmetry, transitivity, and term congruence therefore imply
$r=h\langle t\rangle$. Consequently Refl and Step are interderivable,
with arbitrary side assumptions, with the normalized successor
computation and the old one-template equality rule. This is an exact
consequence-equivalence theorem, not merely a standard-model soundness
observation.

\subsection{A normal zero witness for the single introduction}

The familiar single-introduction convention takes progress at a literal
equation with literal witness $0$, and cases otherwise. Its old alternate
zero witness $I(0,0;u.0)$ is no longer available: it is definitionally
zero in the canonical grammar. Instead use
\[
 q:=I(S0,0;u.u).
\]
This is a closed canonical term which is not the literal $0$.
Successor computation proves $q=0$; equivalently one Step on the
reflexivities of $0$ and $q$ proves it directly. This syntactic distinction
is the only property of the alternate witness needed below.

\section{A five-pattern clausal calculus}
\label{sec:five-pattern-clauses}

This subsection gives the full logical interface. It is independent of
whether a checker physically stores named variables, de Bruijn indices,
or an equivalent scoped encoding.

\subsection{Formulas and retained clauses}

Formulas have only the constructors
\[
 A ::= (s=t)\mid D_{x,y}B(x,y).
\]
The variables $x,y$ are bound in $B$. The ordinary interpretation is
\[
 \bigl(D_{x,y}B\bigr)^*
 :=\exists x\bigl(B(x,0)^*\to\forall y\,B(x,y)^*\bigr).
\]
The designated false formula is $0=S0$.
A retained clause is written
\[
 [A\,;\,y.B(y)]
\]
and has ordinary reading $A^*\to\forall y\,B(y)^*$; its variable $y$
is bound in $B$ and does not occur free in $A$. A context is a finite
list of such clauses. Empty contexts are allowed. For a formula $A$,
its ordinary-assumption encoding is
\[
 \operatorname{Assm}(A):=[0=0\,;\,y.A],\qquad y\notin\operatorname{FV}(A).
\]
This is an abbreviation for a clause, not a sixth inference rule.

\subsection{Complete rule inventory}

The following five patterns, and no further logical or arithmetic rules,
are primitive:
\begin{align*}
 \text{Use:}\quad&
 \frac{[A;y.B]\in\Gamma\qquad\Gamma\vdash A}
      {\Gamma\vdash B\langle t\rangle};\\[2mm]
 \text{Refl:}\quad&\Gamma\vdash t=t;\\[2mm]
 \text{Step:}\quad&
 \frac{\Gamma\vdash I(c,a;u.h)=t\qquad
       \Gamma\vdash I(Sc,a;u.h)=r}
      {\Gamma\vdash r=h\langle t\rangle};\\[2mm]
 \text{DE:}\quad&
 \frac{\Gamma\vdash D_{x,y}B\qquad
       \Gamma,[B(b,0);y.B(b,y)]\vdash A}
      {\Gamma\vdash A};\\[2mm]
 \text{DI:}\quad&
 \frac{\Gamma,\operatorname{Assm}(B(t,\eta))
          \vdash B(t,Sz)}
      {\Gamma\vdash D_{x,y}B}.
\end{align*}
For DE, $b$ is fresh for the outer context, the conclusion, and the
passive parameters. For DI, $z$ is fresh for the outer context, the
witness $t$, and all passive parameters; moreover
\[
 \eta=\begin{cases}
 z,& B\text{ is a literal equation and }t\text{ is literally }0,\\
 0,&\text{otherwise}.
 \end{cases}
\]
Every term and formula substitution in these rules is normalized.
Fresh variables used in a premise enlarge the scope of the entire outer
context as well as the formulas. Omitting that weakening is unsound.

All five patterns have one conclusion. Refl is an axiom schema.
Step supplies both computation and equational inference. Use accesses
retained clauses and thus includes the ordinary hypothesis rule:
from $\operatorname{Assm}(A)\in\Gamma$, apply Use to Refl$(0)$.
The only induction condition appears in the syntactic choice of
$\eta$ at DI; no induction for compound matrices is permitted.

\subsection{Structural bookkeeping}

Context weakening and renaming of free variables are admissible by
induction on derivations. Renaming preserves canonicality, literal zero,
and the root constructor of a formula, so the DI mode is preserved.
Ordinary formula cut is also admissible: if
$\Gamma\vdash A$ and $\Gamma,\operatorname{Assm}(A)\vdash C$, then
$\Gamma\vdash C$. Replace each Use of that assumption clause by the
first derivation, weakened to the current eigenvariable scope; induction
on the second derivation completes the argument. Its antecedent proof
of $0=0$ can be discarded. Bound eigenvariables are renamed before the
replacement when necessary.

These arguments do not require a direct proof of arbitrary term
substitution for the syntax-directed DI. Such a substitution may change
a nonliteral witness into literal zero. General term substitution becomes
admissible after equivalence with the ordinary native kernel has been
proved; treating its preservation of the mode as automatic would be an
error.

\subsection{Conditional closures as formulas}

Define the closure abbreviation
\[
 \widehat Q_y B(y):=D_{x,y}B(y),\qquad x\notin\operatorname{FV}(B).
\]
Its ordinary reading is equivalent to $B(0)\to\forall yB(y)$ because
the numerical domain has the closed inhabitant $0$.

\paragraph{Closure elimination.}
From $\Gamma\vdash\widehat Q_yB$ and $\Gamma\vdash B(0)$, apply DE.
Its retained clause is $[B(0);y.B(y)]$ and is independent of the fresh
witness. Use at any target $t$, with the weakened base proof, gives
$B(t)$. Thus QE is derived.

\paragraph{Closing a retained canonical clause.}
From a retained clause $[B(0);y.B(y)]$, derive $\widehat Q_yB$ by DI
with dummy witness $q$. This is the case mode, since $q$ is not literal
zero. Under the local assumption $B(0)$, Use proves $B(Sz)$.
The same argument converts any derivation
$\Gamma,\operatorname{Assm}(B(0))\vdash B(Sz)$, with $z$ fresh,
into a derivation of $\widehat Q_yB$: take the displayed derivation
as the premise of DI at dummy witness $q$.

\paragraph{Equation progress.}
Given
\[
 \Gamma,\operatorname{Assm}(E(z))\vdash E(Sz),
 \qquad E\text{ a literal equation},
\]
apply DI to the dummy-witness matrix $E(y)$ at witness $0$.
It produces $\widehat Q_yE$. Thus the native progress rule is derived
with no equality reasoning or additional induction schema.

\subsection{Recovering every case introduction}

For a compound matrix, or for a witness other than literal zero, DI is
already the desired case rule. Only a literal equation $E(x,y)$ with
witness $0$ requires an argument. Suppose
\[
 \Gamma,\operatorname{Assm}(E(0,0))\vdash E(0,Sz).
\]
By closing the displayed derivation at the dummy witness $q$, obtain
$\widehat Q_y E(0,y)$. Apply DE to retain
$[E(0,0);y.E(0,y)]$ and introduce $D_{x,y}E(x,y)$ by DI at witness
$q$, which uses the case mode. In its premise assume $E(q,0)$.
Transport along $q=0$ gives $E(0,0)$; Use of the retained clause gives
$E(0,Sz)$; transport back along $0=q$ gives $E(q,Sz)$, as required.
The fresh eigenvariable from the outer DE does not occur in the final
conclusion. All transports are in literal equation templates and are
available from the arithmetic Pivot construction.

Hence the five-pattern calculus derives both original introduction
rules, including the otherwise missing atomic case at zero. This proof
also explains why normalizing the old nonliteral zero witness without
replacing it is insufficient.

\section{Exact equivalence and the inherited iPRA bridge}
\label{sec:uniform-step-equivalence}

Let $N_0$ denote the original native calculus on raw scalar terms, and
let $N_z$ be its zero-normal presentation, with ordinary Hyp, successor
computation, one-template transport, QE, DE, case DI, and atomic QI.
The pinned baseline proves
\[
 N_0:\Gamma\vdash A
 \quad\Longleftrightarrow\quad
 N_z:N\Gamma\vdash NA.
\]
Its formal source is \texttt{NativeZeroNormal.lean}; the normal-form
coherence and interpretation facts are in
\texttt{ZeroNormalization.lean}. The semantic fact that a raw term
is provably equal to its normal form under an iterator body uses
body congruence, which in that inherited development uses equation
induction. The new equality bootstrap above does not claim to avoid
that inherited requirement in the proof of the normalization theorem.

Define $T$ on the formulas of $N_z$ by preserving equality and $D$,
and replacing every $Q_yB$ by $\widehat Q_y(TB)$. Translate ordinary
hypotheses $A$ to $\operatorname{Assm}(TA)$. The translation $T$
commutes with normalized numerical substitution up to renaming bound
variables: prove this by structural induction on formulas, choosing
each dummy witness binder fresh for the substitution images.
The preceding derivations
translate every rule of $N_z$ into the five-pattern calculus $C$:
Hyp becomes Use with Refl$(0)$; successor computation and transport
are derived from Step; QE uses DE and Use; equation progress is DI
at a dummy zero witness; every case introduction is recovered as
above. In the translation of DE, the retained canonical clause derives
the translated hypothetical closure, so admissible formula cut supplies
the translated branch derivation.

For the reverse direction, interpret retained clauses as
$A^*\to\forall yB(y)^*$ in the ordinary scalar theory $T_{\rm eq}$.
Use is implication elimination followed by universal elimination.
Refl is equality reflexivity. Step follows from successor computation
and equality. DE is ordinary existential elimination. In case-mode DI,
coverage $n=0\lor\exists z\,n=Sz$ combines the base assumption and
the successor premise to prove the conditional closure. In progress-mode
DI, contextual induction for the literal equation supplies that closure.
Existential introduction then supplies the stated witness in each mode.
Thus the entire calculus is sound in $T_{\rm eq}$ with precisely the
same domain and actual equality.

The pinned ordinary/native equivalence and the rule translations yield,
for every finite list $\Gamma$ of native formulas and every native $A$,
\[
 N_0:\Gamma\vdash A
 \quad\Longleftrightarrow\quad
 C:\operatorname{Assm}(T(N\Gamma))\vdash T(NA).
 \tag{Faithfulness}
\]
Here the right-hand context means pointwise application to its entries.
The reverse implication uses that the ordinary reading of
$\operatorname{Assm}(B)$ is equivalent to $B$, and that the ordinary
reading of the dummy-existential closure is equivalent to $Q$.
Those equivalences are intuitionistically valid over an inhabited domain.

\paragraph{Bi-interpretability, rather than only equiconsistency.}
The inherited bridge interprets raw scalar terms as their fixed primitive
recursive definitions in iPRA, and compiles each fixed primitive recursive
definition back to a scalar term. Its round trips are provably the identity
on numerical values. Compose those maps with $N$ and $T$. The normalizing
map fixes variables and zero, preserves numerical values provably, and
leaves the whole numerical domain unchanged. The closure replacement has
no numerical effect. Consequently both composites retain the identity
graph $x=y$; equality remains actual equality, and no quotient or proper
subdomain is introduced.

On formulas, inclusion followed by $T\circ N$ is the identity on the
canonical Eq/$D$ language. The other round trip is provably equivalent
to the input formula by zero-normalization equivalence and
$Q_yB\leftrightarrow D_{x,y}B$ for a dummy $x$. Transporting the inherited
ordinary/iPRA formula round-trip theorems through Faithfulness gives both
full first-order round trips for $C$. In particular, the designated
contradiction $0=S0$ is preserved, so the resulting theories are
equiconsistent as well as bi-interpretable.

\paragraph{Boundary of the claim.}
The new mathematical deletion is the replacement of the multi-case
arithmetic equality/computation mechanism by Refl and one uniform Step.
Zero-normalization and clausal logical contexts are inherited design
ideas and their implementation costs must be counted. A numerical size
comparison is a separate claim about complete programs, not a theorem
that no smaller presentation exists. Any proof-assistant status must be
reported according to the files actually compiled; the written
five-pattern equivalence is not automatically a machine-checked theorem.
